%% kapa: exam=1 round=27 year=2016 period=2 form=A subject=accounting subject_ko=회계학 seq=10 number=50
%% source: Q-Net 감정평가사 기출문제 (cst00309) artlSeq=5169402 file=extracted/1008/5169402/제27회 감정평가사 1차 문제/제27회 감정평가사 2교시문제(A형).pdf
%% conversion: kordoc 4.24.51 -> tools/gen1.py (auto); review=proofread-sonnet; flags=none
\begin{question}{50}
(주)감평은 20×1년 1월 1일에 종업원 100명에게 각각 10개의 주식선택권을 부여하고 4년의 용역제공조건을 부과하였다. 부여시점의 주식선택권 공정가치는 개당 ￦10이다. (주)감평은 종업원 중 20명이 부여일로부터 4년 이내에 퇴사하여 주식선택권을 상실할 것으로 추정하였으나 20×1년 말까지 실제로 퇴사한 종업원은 없었다. 20×2년 말에는 가득기간 동안 30명이 퇴사할 것으로 추정을 변경하였으며 20×2년 말까지 실제 퇴사한 종업원은 없었다. 주식선택권의 부여와 관련하여 20×2년도에 인식할 보상비용은?
\begin{choices}
  \item ￦1,000
  \item ￦1,500
  \item ￦1,750
  \item ￦2,000
  \item ￦2,500
\end{choices}
\end{question}
\answer{2}
